%\input ../util/bookmacros
%\input ../util/docparams

\appendix{Appendix E}{Taylor's Theorem}

In the chapter ``The Derivative As Approximation'' we approximated a function near a point by its
tangent line -- $f(x+h) \approx f(x) + f'(x) h$ -- and we were rather vague about the size of the
error we were making. That chapter left two questions hanging: exactly how big is the error, and
can we shrink it by adding a quadratic correction, or a cubic one, and so on? The answer to both
questions is {\it Taylor's Theorem\/}, which we prove carefully in this appendix. Its simplest
cases return exactly the tangent line approximation -- but now with a formula, not a hand wave,
for the error.

\section{Generalization of the Mean Value Theorem}
We start by stating and proving a generalization of the Mean Value Theorem:
\proclaim Theorem(\index{Generalized Mean Value Theorem}). Let $f$ and $w$ be continuous functions 
on the interval $[a,b]$. If $f$ has minimum and maximum values $m$ and $M$ while 
$w(s)  \ge 0$ for $s \in [a,b]$ and $\int_a^b w(s) \, ds > 0$,  then 
there exists a value $\mu \in [a,b]$ such that
$$
\eqalign{
\int_a^b f(s) w(s)\, ds & = f(\mu) \int_a^b w(s) \, ds \cr
}
$$

{\bf proof:\/} Since $f$ is greater or equal to $m$ and less than or equal $M$ we must have
$$
\int_a^b m\, w(s) \, ds \le \int_a^b f(s) w(s) \, ds \le \int_a^b M w(s) \, ds
$$
Since $m$ and $M$ are constants they may be taken out of the integrals (linearity property of integrals), 
so that
$$
m\int_a^b w(s) \, ds \le \int_a^b f(s) w(s) \, ds \le M \int_a^b w(s) \, ds
$$
Dividing by the integral in $w$ gives
$$
m \le {\int_a^b f(s) w(s) \, ds \over \int_a^b w(s) \, ds} \le M 
$$
Note that the middle term is an intermediate value of the function $f$.
Since $f$ is a continuous function on $[a,b]$, the \index{intermediate value theorem} tells us that 
there is a $\mu\in[a,b]$
such that $f(\mu) = {\int_a^b f(s) w(s) \, ds \over \int_a^b w(s) \, ds}$. This can be rewritten as
$$
\int_a^b f(s) w(s) \, ds = f(\mu) \int_a^b w(s) \, ds
$$

\proclaim Corollary(Mean Value Theorem). If $f$ is a function such that $f^{\prime}$ exists and is a continuous function
on the interval $[a,b]$, then there exists a $\mu \in [a,b]$\footnote{\kern 0.5pt \raise 0.5ex \hbox{\dag}}{\rm The
Mean Value Theorem is usually stated with weaker hypotheses -- $f$ need only be differentiable on the {\it open\/}
interval $(a,b)$ -- and with a sharper conclusion -- $\mu$ lies strictly between $a$ and $b$. Proving that version
requires other tools; this integral version is all we need here.} such that
$$ 
\eqalign{
f(b) & = f(a) + f'(\mu)(b - a) \cr
}
$$

This follows from the generalized Mean Value result substituting $f'$ for $f$ and setting $w(x) \equiv 1$ for all $x \in [a,b]$.
That is, we have for some $\mu \in [a,b]$:
$$
\eqalign{
\int_a^b f^\prime(s) 1 \, ds & = f^\prime(\mu) \int_a^b 1 \, ds \cr
}
$$
Or,
$$
  f(b) - f(a) = f^\prime(\mu) (b - a)
$$
\section{Proof of Taylor's Theorem}
In the chapter ``The Derivative As Approximation'' we derived Taylor's Theorem by pretending that
the various ``$\mu$'' values handed to us by the Mean Value Theorem were all the same -- and we
confessed there that this was a cheat. Here we take an honest route. The only tools we need are
the \index{FTC} and \index{Integration by Parts} -- both from the chapter on integral calculus --
together with a {\it single\/} application of the Generalized Mean Value Theorem at the very end.

Suppose $f'$ is continuous on the interval $[a,b]$ and let $x \in [a,b]$. The \index{FTC} tells us that
$$
\eqalignno{
f(x) & = f(a) + \int_a^x f'(s) \, ds & (1) \cr
}
$$
Squint at this equation and you can already see the shape of Taylor's Theorem: a polynomial
approximation to $f$ (here just the constant $f(a)$) plus an integral that records the error
{\it exactly\/}.

To do better we apply \index{Integration by Parts} to the error integral, using the two functions
$f'(s)$ and $v(s) = -(x - s)$. Keep in mind that we are integrating with respect to $s$, so the
value $x$ is a constant as far as the integration is concerned; therefore, $v'(s) = 1$, and we may
write $\int_a^x f'(s)\, ds = \int_a^x f'(s)\, v'(s)\, ds$.
Why the minus sign in $v$? An anti-derivative of $1$ is only determined up to a constant, and this
choice makes $v$ vanish at the endpoint $s = x$ -- watch what that does for us.
If $f''$ is also continuous on $[a,b]$, then Integration by Parts gives
$$
\eqalign{
\int_a^x f'(s)\, v'(s)\, ds & = \bigl(f'(x)\, v(x) - f'(a)\, v(a)\bigr) - \int_a^x f''(s)\, v(s)\, ds \cr
& = f'(a) (x - a) + \int_a^x f''(s) (x - s)\, ds \cr
}
$$
The endpoint terms collapsed to $f'(a)(x-a)$ because $v(x) = -(x - x) = 0$ while $-v(a) = (x - a)$;
and the minus sign in front of the last integral canceled the minus sign inside $v$.
Plugging this back into (1) gives
$$
\eqalignno{
f(x) & = f(a) + f'(a) (x - a) + \int_a^x f''(s) (x - s)\, ds & (2) \cr
}
$$
This is the tangent line approximation from the chapter ``The Derivative As Approximation'' --
but now with an exact formula for the error.

Let's do it once more, applying Integration by Parts to the error integral in (2), this time with
the two functions $f''(s)$ and $v(s) = -{(x - s)^2 \over 2}$. By the chain rule $v'(s) = (x - s)$ --
differentiating $(x - s)$ with respect to $s$ produces a factor of $-1$ that cancels the minus sign
out front -- and once again $v(x) = 0$. If $f'''$ is also continuous on $[a,b]$, then
$$
\eqalign{
\int_a^x f''(s) (x - s)\, ds & = \bigl(f''(x)\, v(x) - f''(a)\, v(a)\bigr) - \int_a^x f'''(s)\, v(s)\, ds \cr
& = f''(a) {(x - a)^2 \over 2} + \int_a^x f'''(s) {(x - s)^2 \over 2}\, ds \cr
}
$$
so that (2) becomes
$$
\eqalignno{
f(x) & = f(a) + f'(a) (x - a) + f''(a) {(x - a)^2 \over 2!} + \int_a^x f'''(s) {(x - s)^2 \over 2!}\, ds & (3) \cr
}
$$
Each round of Integration by Parts peels one more polynomial term off of the error integral and
leaves behind an error integral of the same shape. The pattern that emerges is this: if
$f^{(k+1)}$ exists and is continuous on $[a,b]$, then
$$
\eqalignno{
f(x) & = f(a) + \sum_{i=1}^{k} f^{(i)}(a) {(x - a)^i \over i!} + \int_a^x f^{(k+1)}(s) {(x - s)^k \over k!}\, ds & (4) \cr
}
$$
Call this statement $P_k$. Notice that the hypothesis is honest bookkeeping: each round of
Integration by Parts differentiates $f$ once more under the integral, so reaching $P_k$ costs us
$k+1$ continuous derivatives -- no hidden assumptions about mysterious ``$\mu$'' functions.

How do we prove $P_k$ for every $k$? We use \index{induction} (see the appendix on induction):
we must show that $P_0$ is true and that $P_k$ implies $P_{k+1}$ for any $k \ge 0$.

$P_0$ is true since it is just equation (1) -- when $k = 0$ the sum is empty while
${(x - s)^0 \over 0!} = 1$.

For the inductive step, assume that $P_k$ is true and that $f^{(k+2)}$ exists and is continuous
on $[a,b]$. We apply Integration by Parts to the error integral in (4) exactly as before, this
time with the two functions $f^{(k+1)}(s)$ and $v(s) = -{(x - s)^{k+1} \over (k+1)!}$. The power
rule and chain rule give
$$
v'(s) = {(k+1) (x - s)^k \over (k+1)!} = {(x - s)^k \over k!}
$$
-- as usual the two minus signs cancel -- and, as usual, $v(x) = 0$. Therefore
$$
\eqalign{
\int_a^x f^{(k+1)}(s) {(x - s)^k \over k!}\, ds & = \bigl(f^{(k+1)}(x)\, v(x) - f^{(k+1)}(a)\, v(a)\bigr) - \int_a^x f^{(k+2)}(s)\, v(s)\, ds \cr
& = f^{(k+1)}(a) {(x - a)^{k+1} \over (k+1)!} + \int_a^x f^{(k+2)}(s) {(x - s)^{k+1} \over (k+1)!}\, ds \cr
}
$$
Substituting this into (4) extends the sum to $i = k+1$ and produces exactly the statement
$P_{k+1}$. The two requirements for \index{induction} have been fulfilled; therefore (4) is true
for every $k \ge 0$.

One task remains. The error integral in (4) is exact, but it is not easy to look at. We convert it
into a friendlier form by applying the \index{Generalized Mean Value Theorem} -- just once -- on the
interval $[a,x]$, with $f^{(k+1)}$ playing the role of $f$ and with the weight function
$w(s) = {(x - s)^k \over k!}$. (For $x = a$ both sides of (4) are $f(a)$, so there is nothing to
prove; since $x \in [a,b]$, we may assume $x > a$.) Let us check the hypotheses: $f^{(k+1)}$
and $w$ are continuous, and $w(s) \ge 0$ for $s \in [a,x]$ since $s \le x$ there. Also, an
\index{anti-derivative} of ${(x - s)^k \over k!}$ with respect to $s$ is
$-{(x - s)^{k+1} \over (k+1)!}$ -- our friend $v$ again -- so by the \index{FTC}
$$
\int_a^x {(x - s)^k \over k!}\, ds = -{(x - x)^{k+1} \over (k+1)!} + {(x - a)^{k+1} \over (k+1)!} = {(x - a)^{k+1} \over (k+1)!}
$$
which is positive since $x > a$. The Generalized Mean Value Theorem then hands us a value
$\mu \in [a,x]$ -- a value that depends on $x$, so it is really a function $\mu(x)$ -- such that
$$
\eqalign{
\int_a^x f^{(k+1)}(s) {(x - s)^k \over k!}\, ds & = f^{(k+1)}(\mu(x)) \int_a^x {(x - s)^k \over k!}\, ds = f^{(k+1)}(\mu(x)) {(x - a)^{k+1} \over (k+1)!} \cr
}
$$
Substituting this into (4), and then writing $n$ for $k + 1$, turns (4) into the celebrated result:
\proclaim Theorem(Taylor's Theorem). If $f$ is a function such that $f^{(n)}$ exists and is a continuous function 
on the interval $[a,b]$, then for any $x \in [a,b]$\footnote{\kern 0.5pt \raise 0.5ex \hbox{\dag}}{\rm Note, when $n=1$,
  $\sum_{i=1}^{n-1} g_i = \sum_{i=1}^0 g_i = 0$, regardless of the values of $g_i$. }
$$
\eqalign{
f(x) &= f(a) + \sum_{i=1}^{n-1} f^{(i)}(a) {(x - a)^i \over i!} + f^{(n)}(\mu(x)) {(x-a)^n \over n!} \cr
}
$$
where $\mu$ is a function such that $\mu(x) \in [a,x]$.

Notice that when $n = 1$ the theorem says $f(x) = f(a) + f'(\mu(x)) (x - a)$ -- the Mean Value
Theorem with which we began, now seen as the first rung of a ladder.

Let us make the error term earn its keep.

\example{How good is the approximation $\sin(x) \approx x$ when $x$ is small?}
With $f(x) = \sin(x)$ and $a = 0$ we have $f(0) = 0$, $f'(0) = \cos(0) = 1$, and
$f''(0) = -\sin(0) = 0$ -- so the tangent line at $0$ is the line $y = x$, and the quadratic
correction vanishes for free. Taylor's Theorem with $n = 3$ then says
$$
\eqalign{
\sin(x) & = x + f'''(\mu(x)) {x^3 \over 3!} \cr
}
$$
for some $\mu(x)$ between $0$ and $x$. Now $f''' = -\cos$, which is never bigger than $1$ in
size -- no matter what $\mu(x)$ turns out to be. So for $|x| \le 0.1$,
$$
| \sin(x) - x | \le {|x|^3 \over 3!} \le {(0.1)^3 \over 6} \approx 1.7 \times 10^{-4}
$$
Not bad for a straight line -- and we never had to compute $\mu(x)$.

\vfil
\eject

